\chapter*{Agradecimientos}

Hay personas sin las cuales este trabajo no habría sido posible, y hay otras sin las cuales tampoco habría valido la pena hacerlo. A todas ellas van dedicadas estas líneas.

A mi padre, por estar siempre presente. Por cada vez que me preguntó cómo iba el TFG sin entender del todo la respuesta, pero escuchándome igual con toda su atención. Gracias por creer en mí incluso cuando yo no lo hacía.

A \textbf{Iana Miranda} y \textbf{Macarena Pereira}, por haber sido mucho más que compañeras de carrera. Por las tardes de estudio, las risas entre exámenes, el apoyo en los momentos difíciles y la alegría en los buenos. Sois de las personas más bonitas que me ha dado esta etapa, y espero que sigáis siendo parte de las que están por venir.

A \textbf{Irene González} e \textbf{Ignacio Mora}, por sus aportaciones, su tiempo y su disposición a ayudar a mejorar este proyecto. Gracias por tomarse el interés que se tomaron, por las ideas que aportaron y por hacerme ver cosas que yo sola no habría visto.

A \textbf{Santia}, porque ha sido la persona con quien he vivido este proceso de principio a fin. Cuatro años de carrera juntas, pero el TFG ha sido algo diferente: las dudas compartidas, los agobios que nos entendíamos sin tener que explicarlos demasiado, las victorias pequeñas que solo alguien que está pasando por lo mismo puede entender de verdad. Gracias por estar ahí en cada una de esas fases, por la paciencia, por el apoyo constante y por hacer que este camino se sintiera menos solitario.

Y por último, a \textbf{José Antonio Parejo}, por ser exactamente el tutor que necesitaba. Por su implicación, su cercanía, su paciencia y por crear siempre un ambiente de trabajo en el que me sentí escuchada y acompañada. Este trabajo también es suyo.
